\setlength{\absparsep}{18pt} % ajusta o espaçamento dos parágrafos do resumo
\begin{resumo}[Abstract]

	\begin{otherlanguage*}{english}

		This work presents the development of a web platform for creating and sharing digital collections of any kind, motivated by the rigidity of the usual record-keeping tools, which treat each record as a set of fixed fields with a standardized appearance. In the platform, the user's holdings are organized into categories, collections and items, and the card of each item is defined by the user in a visual layout editor with six types of elements. The proposal is completed by the customization of the interface through color palettes and typography and by social features of friendship, collection following, feed and notifications, under the privacy set for each collection. The development started from the specification of the requirements and the domain model, followed a layered architecture, with a React interface, a NestJS API and a PostgreSQL database, and was verified through automated tests, manual verification and performance measurements. The platform has been published and is in operation. All 546 automated tests, unit and end-to-end, passed, the twenty-four functional requirements were met, and the performance metrics were met by a wide margin: with the network limited to 10 Mbps, the measured screens displayed their main content in less than half a second at the 75th percentile, and the editor responded to interactions within 48 milliseconds. It is concluded that the platform brings together, in a single place, the organization of collections, the freedom to define the structure and the appearance of each record, and sharing with other people.

		\vspace{\onelineskip}

		\textbf{Keywords}: Digital collections. Interface customization. Layout editor. Web application. Sharing.

	\end{otherlanguage*}

\end{resumo}
